%% --- zde jsou zavedeny některé "konstanty" - některé musíte změnit! --- %%
\newcommand{\cvut}{České vysoké učení technické v~Praze}
\newcommand{\fjfi}{Fakulta jaderná a fyzikálně inženýrská}
\newcommand{\ksi}{Katedra softwarového inženýrství}
\newcommand{\km}{Katedra matematiky}
% \newcommand{\program}{Aplikace přírodních věd} % změňte, pokud máte jiný stud. program
\newcommand{\obor}{Aplikace informatiky v přírodních vědách} % změňte, pokud máte kurzívujiný obor

\newcommand{\druh}{Výzkumný úkol} % nebo "Diplomová práce"
\newcommand{\woman}{a} % pokud jste ŽENA, ZMĚŇTE na: ...{\woman}{a} (je to do Prohlášení)

\newcommand{\logoCVUT}{\includegraphics{symbol_cvut_konturova_verze_cb.pdf}} % logo ČVUT -- podle grafického manuálu ČVUT platného od prosince 2016. Pokud nevyhovuje PDF-verze, tak použijte jinou variantu loga: https://www.cvut.cz/logo-a-graficky-manual -> "Symbol a logo ČVUT v Praze"). Pokud chcete logo úplně vynechat, zadejte místo "\includegraphics{...}" text "\vspace{35mm}"

% přesně podle formuláře "Zadání bak./dipl. práce" VYPLŇTE:
\newcommand{\nazevczT}{Optické rozpoznávání znaků \\ na naskenovaných historických plakátech pomocí \\nejmodernějších metod}    % český název práce (přesně podle zadání!)
\newcommand{\nazevenT}{Optical Character Recognition \\on Scanned Historical Posters Using the State-of-the-Art Methods}          % anglický název práce (přesně podle zadání!)
\newcommand{\nazevcz}{Optické rozpoznávání znaků na naskenovaných \\historických plakátech pomocí nejmodernějších metod}    % český název práce (přesně podle zadání!)
\newcommand{\nazeven}{Optical Character Recognition on Scanned Historical Posters Using the State-of-the-Art Methods}          % anglický název práce (přesně podle zadání!)
\newcommand{\autor}{Anna Gruberová}   % vyplňte své jméno a příjmení (s akademickým titulem, máte-li jej)
\newcommand{\vedouci}{Ing. Adam Novozámský, Ph.D.} % vyplňte jméno a příjmení vedoucího práce, včetně titulů, např.: Doc. Ing. Ivo Malý, Ph.D.
\newcommand{\pracovisteVed}{Computer Vision Lab, Institute of Visual Computing \& Human-Centered Technology, TU Wien - Faculty of Informatics} % ZMĚŇTE, pokud vedoucí Vaší práce není z KSI
\newcommand{\konzultant}{--} % POKUD MÁTE určeného konzultanta, NAPIŠTE jeho jméno a příjmení
\newcommand{\pracovisteKonz}{--} % POKUD MÁTE konzultanta, NAPIŠTE jeho pracoviště

% podle skutečnosti VYPLŇTE:
\newcommand{\rok}{2022}  % rok odevzdání práce (jen rok odevzdání, nikoli celý akademický rok!)
\newcommand{\kde}{Praze} % studenti z Děčína ZMĚNÍ na: "Děčíně" (doplní se k "prohlášení")

\newcommand{\klicova}{Optické rozpoznávání znaků, rozpoznávání textu, automatizace}   % zde NAPIŠTE česky max. 5 klíčových slov
\newcommand{\keyword}{Optical Character Recognition, Text Recognition, Automation}       % zde NAPIŠTE anglicky max. 5 klíčových slov (přeložte z češtiny)
\newcommand{\abstrCZ}{Optické rozpoznávání znaků z obrazových dat je žádanou úlohou v dnešním světě, jelikož pro člověka je již nemožné bez automatizace zpracovat velké množství obrazových dat. Vídeňská knihovna vlastní nad 350 tisíc digitalizovaných historických plakátů, z nichž je třeba extrahovat zobrazený text. Cílem této práce je zmapovat vybrané existující metody rozpoznávání textu a porovnat je na základě testů na vybraných datasetech.}


% zde NAPIŠTE abstrakt v češtině (cca 7 vět, min. 80 slov)
\newcommand{\abstrEN}{Optical character recognition from image data is a demanding task in today's world, because it is already impossible for humans to process a large amount of image data without automation. The Vienna Library owns over 350,000 digitized historical posters, from which the displayed text must be extracted. The aim of this work is to map selected existing text recognition methods and compare them based on tests on selected datasets.} % zde NAPIŠTE abstrakt v angličtině

\newcommand{\prohlaseni}{Prohlašuji, že jsem svou bakalářskou práci vypracoval\woman{} samostatně a použil\woman{} jsem pouze podklady (literaturu, projekty, SW atd.) uvedené v přiloženém seznamu.} % text prohlášení můžete mírně upravit :-)

\newcommand{\podekovani}{.} % NAPIŠTE poděkování, např. svému vedoucímu:
% Děkuji Ing. Eleonoře Krtečkové, Ph.D. za vedení mé bakalářské práce a za podnětné návrhy, které ji obohatily.
% NEBO:
% Děkuji vedoucímu práce doc. Pafnutijovi Snědldítětikaši, Ph.D. za neocenitelné rady a pomoc při tvorbě bakalářské práce.

\newcommand{\ti}{\textit} % zkrácený příkaz pro kurzívu
\newcommand{\tb}{\textbf} % zkrácený příkaz pro tučné písmo
\newcommand{\tn}{\texttt} % zkrácený příkaz pro neproporcionalni písmo

% Vzhled kodu - lstlistings - python
\definecolor{codegreen}{rgb}{0,0.6,0}
\definecolor{codegray}{rgb}{0.5,0.5,0.5}
\definecolor{framegray}{rgb}{0.8,0.8,0.8}
\definecolor{codepurple}{rgb}{0.58,0,0.82}
\definecolor{backcolour}{rgb}{0.95,0.95,0.92}

\lstdefinestyle{mystyle}{
    language=Python,
    frame=single,
    rulecolor=\color{framegray},
    commentstyle=\color{codegreen},
    keywordstyle=\color{magenta}\bfseries,     numberstyle=\tiny\color{codegray},
    stringstyle=\color{codepurple},
    basicstyle=\ttfamily\footnotesize,
    breakatwhitespace=false,         
    breaklines=true,                 
    captionpos=t,                    
    keepspaces=true,                 
    numbers=left,                    
    numbersep=5pt,                  
    showspaces=false,                
    showstringspaces=false,
    showtabs=false,                  
    tabsize=2
}
\lstset{style=mystyle}

\renewcommand{\lstlistingname}{Function}

%-----------------------------------------------

\newenvironment{repl}
{\fontfamily{cmvtt}\selectfont \begin{mdframed}

}{\end{mdframed}}

\usepackage[linewidth=1pt]{mdframed}


\newcolumntype{L}{>{\centering\arraybackslash}m{2cm}}
\newcolumntype{N}{>{\centering\arraybackslash}m{1.7cm}}
\newcolumntype{M}{>{\centering\arraybackslash}m{1.6cm}}
\newcolumntype{K}{>{\centering\arraybackslash}m{1.5cm}}
\newcolumntype{S}{>{\centering\arraybackslash}m{1cm}}
\renewcommand{\arraystretch}{1.3}

\newcommand{\cmark}{\textcolor{green!80!black}{\ding{51}}}
\newcommand{\xmark}{\textcolor{red}{\ding{55}}}


